% !TeX root = ../../surface-lagrangian-torus-fibration.tex
% 2026-11-21 / v1.2: Component の対数費用と残余基底の先行選択による全域上界．
\section{全 fiber 合併の収縮と actual 全曲面距離の上界}
\label{ke:global-ambient-upper-section}

命題\ref{ke:global-base-parameterization}で構成した
$f_{j,t}=p_T\pi_{j,t}:X_t\to\mathcal C_d$ を用いる．
本節の全ての source 側の距離は，\emph{全曲面の経路を許す actual ambient 距離}
$d_{g_t}$ である．保持 corner 自身の内在距離は用いない．
十分大きい固定 stage $j$ の後に $|t|$ を小さくし，記号を
\[
 \bar d_t=R_t^{-1}d_{g_t},\qquad m=m_j=6b_j,\qquad
 \delta=4\pi/N_j,\qquad
 \eta=\eta_{j,t}=A\frac{1+\log(N_j+1)}{R_t},\qquad\kappa=\eta
\]
とする．旧 corner theorem の誤差を
\begin{equation}\label{ke:eq:global-upper-corner-error}
 E=E_{j,t}=C\left(\frac1{m_j}
       +\frac{\log(m_j+1)+\log R_t}{R_t}\right)
\end{equation}
と書く．以下では $C$ は固定した次数・退化族・有限 chart 系のみに依存し，
有限定数 $C_j$ は stage $j$ に依存してよい．

\subsection{同一 moment の接続費用と全 fiber 合併}

各 pair $ab$ について，neck の全 marked tree と全 source discs を含む連結な
固定残余基底 $P_{ab,j}$ を取る．Source discs を除いた部分も連結に取れる．
さらに，各 labelled end の旧 softmax boundary $w_c=0$ の residual image を
その辺全体について含むようにする．\eqref{ke:eq:hexagon-edge-profile}より
$B=m_j+r_c$，$0\le r_c<1$ なので，この選択は $N_j$ と
$[-\delta,\delta]$ を決める前に行える．既存 K2 の trim $B=7b_j$ より内側にあり，
既に用意した広い fixed-stage residual domain に
$P_{ab,j}\Subset P^+_{ab,j}$ を選べる．
相対葉族補間もこの固定基底の境界 collar を保つので，後から基底を拡大し直さない．

$C_j^{\rm neck}$ を，系\ref{ke:source-inclusive-same-moment-collapse}の
この有限個の固定基底に対する接続定数の最大とする．これは $N_j$ より先に固定できる．
成分側に残る短い neck 尾部にも，次の深度評価がある．

\begin{lemma}[Component 側の neck 尾部の深度]
\label{ke:component-neck-tail-depth}
$N_j$ の下限を大きくした後に $t$ を小さくすれば，共通 anchor interval の外側にある
各 neck 尾部は，対応する一つの component graph 上の
$\max_as_a\le CN_j$ に含まれる．$C$ は stage に依存しない．
従って尾部と同じ component の core・変更 collar・link の任意の二点の接続費用は
$C(1+\log(N_j+1))$ 以下である．
\end{lemma}

\begin{proof}
Component に向く正の moment を $q=\pm\widehat\mu_N$ とすると，尾部は
$q\ge\delta=4\pi/N_j$ の側にある．
補題\ref{ke:component-meridian-normalization}，\ref{ke:closed-component-link}の
fixed-end 評価を，固定した $P^+_{ab,j}$ の全 phase に適用する．
裸の normal 深度 $s=2N_j$ では
\[
 q=\frac{2\pi}{N_j}+o_j(N_j^{-1})+o_t(1)<\frac{4\pi}{N_j}
\]
となるよう，先に $N_j$，最後に $t$ を選べる．Normal 動径に沿う単調性から，
$q\ge\delta$ の component 側の尾部は $s\le2N_j$ に入る．
Source 軸でも同じ向きの正則 normal coordinate を使う．

この評価を同じ actual normal Moser map の像へ移す．
まず $N_j$ より先に固定深度 $M_j$ を取り，固定基底の定数について
$C_{P^+_j}/M_j<\epsilon_0$ を課す．この長さは残余方向の固定 buffer の半分よりも
小さくする．$2M_j\le s\le2N_j$ では，反対側の normal depth は
$L_t-s-O_j(1)\ge L_t-2N_j-O_j(1)$ なので，最後に $t$ を小さくして
これも $2M_j$ 以上にする．Source 軸では反対側の深度は無限大である．従って
補題\ref{ke:fixed-stage-stratum-coverage}の flow 制限と ODE 一意性により，
元と同じ centered flow の軌跡長を $C_{P^+_j}/M_j$ 以下にできる．
Fixed frame 差と残余深度 $O(b_j)$ を $N_j$ の下限に吸収すると，
初期点の全 component depths は $3N_j$ 以下である．
固定 cusp 比較により，stage に依存しない $K$ で
$|d\log s_a|_{g_{\infty,i}}\le K$ である．
$\max s_a\le6N_j$ の拡大 graph compact 上で
$E_{i,t}^*g_t\ge g_{\infty,i}/2$ を課し，
$\sqrt2K\epsilon_0<\log2$ とする．
すると first-exit argument により軌跡は深度 $6N_j$ の境界へ達しない．
残余方向の退出も，予め確保した $P_{ab,j}\Subset P^+_{ab,j}$ の buffer で防げる．

浅い部分 $s\le2M_j$ の flow 像は，固定 component 上の flow convergence により
先に一つの固定 compact に入る．その全深度の有限上界 $A_j$ を決め，
$N_j\ge A_j+M_j$ を追加する．従って浅い部分にも同じ $CN_j$ 上界がある．
ここまでの右辺の定数は全て $N_j$ より先に固定されている．
最後に命題\ref{ke:component-logarithmic-connectors}を用いれば，接続費用を得る．
Core の深度上界を neck 尾部へ無条件に延長する議論ではない．
\end{proof}

以上から，全接続の予算はまず
$C(1+\log(N_j+1))+C_j^{\rm neck}$ と取れる．
既存の $N_j$ の下限に
\begin{equation}\label{ke:eq:neck-constant-depth-choice}
 C_j^{\rm neck}\le1+\log(N_j+1)
\end{equation}
を加える．これは許容される cut の選択条件であり，$C_j^{\rm neck}$ 自体の
exhaustion に関する増大率を推定した結果ではない．定数 $C$ を一度増やせば
\begin{equation}\label{ke:eq:global-connector-budget}
 H=H_{j,t}:=C\bigl(1+\log(N_j+1)\bigr)
\end{equation}
により次の費用を同時に抑えられる：同一 pair の同一 actual normal moment 上の
任意の二点の接続，各 component の core・変更 collar・neck 尾部内の任意の二点の接続，
および必要な actual normal Moser trajectories．
第一の場合の二点は，異なる reduced 葉又は source・saddle の特異葉に属してよい．
経路自体が moment level 又は一つの fiber 内に留まるとは主張しない．
第二の場合には一つの component ごとに，少し広い固定連結 compact image 内の
経路を使う．異なる component cores 全体の径を $H$ で抑えるわけではない．

\begin{proposition}[潰される木上の全 fiber 合併を含む一様収縮]
\label{ke:collapsed-tree-full-fiber-collapse}
定数 $C$ を十分大きく選べば，全 $z\in\mathcal C_d$ に対し
\begin{equation}\label{ke:eq:full-fiber-diameter-upper}
 \operatorname{diam}_{d_{g_t}} f_{j,t}^{-1}(z)\le H_{j,t}.
\end{equation}
特に，$d\ge6$ で $p_T$ が潰す内部木上の\emph{全ての}
$\pi_{j,t}$-fibers の合併に同じ評価が成立する．
\end{proposition}

\begin{proof}
$\operatorname{supp}z=\{a,b\}$ とし，$z=E_{ab}(u)$ と書く．
$\kappa<u<2-\kappa$ なら，
命題\ref{ke:collapsed-tree-common-moment}により
$f_{j,t}^{-1}(z)$ は共通の
$\mu=\sigma_t^{-1}(u)$ を持つ構成後の neck fibers の合併である．
相対補間後も normal moment は保たれるので，この合併は
$P_{ab,j}$ 上の同一 moment 集合に含まれる．
従って\eqref{ke:eq:global-connector-budget}で径が抑えられる．
両端の $0<u\le\kappa$ 又は $2-\kappa\le u<2$ では，逆像全体が対応する
一つの component の固定 compact core と変更 collar に含まれ，その接続費用を使う．
$z=e_a$ の場合も同じ component compact 内にある．

$|\operatorname{supp}z|=3$ の場合，$p_T$ の逆像は一つの triangle cell の一点で
あり，$f_{j,t}^{-1}(z)$ は一つの構成後の $\pi_{j,t}$-fiber である．
全域構成の pieces とその overlaps は全 fibers で飽和しているので，この fiber は
一つの保持 corner，neck，または component core・変更 collar に含まれる．
保持 corner なら旧 fiber と同じ全 torus であり，
定理\ref{ke:slow-cut-closed-corner-gh}の非正規化 intrinsic diameter $C/m_j$ を使う．
Neck なら同一 moment 集合の部分集合として上の評価を使い，
core・変更 collar なら固定 compact の費用を使う．
$H\ge C/m_j$ となるよう定数を増やせば全ての場合を得る．
\end{proof}

\subsection{共通 moment を持つ corner portals}

三元集合 $J=\{a,b,c\}$ について，旧 softmax 辺 $w_c=0$ 上で
normal moment が $\mu\in[-\delta,\delta]$ である全 torus を
\[
 \Lambda_{ab,J,t}(\mu)\subset A_{t,J}
\]
と書く．補題\ref{ke:common-edge-moment-profile}によりこれは存在し，
\[
 P^{\rm old}_{t,J}(\Lambda_{ab,J,t}(\mu))
       =E_{ab}(u_{J,t}(\mu)).
\]
系\ref{ke:common-neck-longitude}の endpoint rescaling と共通 profile の評価から
\begin{equation}\label{ke:eq:portal-profile-error}
 \sup_{|\mu|\le\delta}|u_{J,t}(\mu)-\sigma_t(\mu)|\le C\eta .
\end{equation}
この評価は全 labelled ends に同じ $\sigma_t$ を用いている．
各 phase の選択は任意でよく，その違いも同一 moment の接続費用に含める．

\begin{lemma}[隣接 corner 間で整合する portals]
\label{ke:coherent-corner-portals}
任意の $z\in\mathcal C_J$ に対し $a_J(z)\in A_{t,J}$ を選び，
\begin{equation}\label{ke:eq:corner-portal-target-error}
 d_\square(P^{\rm old}_{t,J}(a_J(z)),z)\le C\eta
\end{equation}
とできる．さらに，$z\in\mathcal C_J\cap\mathcal C_{J'}$ なら
\begin{equation}\label{ke:eq:corner-portal-junction-cost}
 d_{g_t}(a_J(z),a_{J'}(z))\le H
\end{equation}
となるよう同時に選べる．同じ $J$ を連続して使う場合は，同じ点を使ってよい．
この選択の連続性は主張せず，必要な有限個の gallery endpoints のみに用いる．
\end{lemma}

\begin{proof}
$|\operatorname{supp}z|=3$ なら，この support を含む三元集合は $J$ 自身しかない．
$z\in P^{\rm old}(A_{t,J})$ なら，その旧 fiber から一点を選ぶ．それ以外では
$P^{\rm old}(A_{t,J})$ が $\mathcal C_J$ の $2\eta$-net であることを使って
\eqref{ke:eq:corner-portal-target-error}を満たす一点を選ぶ．
特に $z=P^{\rm old}(x)$ となる指定済みの $x\in A_{t,J}$ があれば，
その $x$ を選んでもよい．異なる corner への接続はこの場合には生じない．

$z=E_{ab}(u)$，$0<u<2$ の場合を考える．
$\kappa\le u\le2-\kappa$ なら
$\mu=\sigma_t^{-1}(u)$ とし，各 $J\supset\{a,b\}$ に対し
$a_J(z)\in\Lambda_{ab,J,t}(\mu)$ を選ぶ．
\eqref{ke:eq:portal-profile-error}が target 誤差を与え，
全ての $a_J(z)$ は同じ $P_{ab,j}$ の同一 moment 集合に属するので
\eqref{ke:eq:global-connector-budget}が接続費用を与える．
$u=1$，すなわち $\mu=0$ でも全 moment 退出定理が適用されるため，
source の pinched fibers があることによる例外はない．

$0<u<\kappa$ なら，component $a$ の link と各保持 corner の交わりから
$a_J(z)\in S_{a,j,t}\cap W_{t,J}$ を一つ選ぶ．
補題\ref{ke:corner-cap-disk}の proper cap arc によりこの交わりは空でなく，
補題\ref{ke:removed-corner-target-cap}により
$P^{\rm old}(a_J(z))$ は $e_a$ の $C\eta$ 近傍にある．
$d_\square(z,e_a)=u<\kappa=\eta$ なので target 誤差を得る．
各 $a_J(z)$ は同じ component compact に属するため，対数接続費用 $H$ を使える．
$2-\kappa<u<2$ も component $b$ を使って同様であり，
$z=e_a$ もこの link 上の選択で扱える．
以上は target の全 support 型を尽くす．
\end{proof}

\begin{lemma}[任意点から gallery の最初の corner への接続]
\label{ke:point-to-corner-portal}
任意の $x\in X_t$ に対し $q_x\in\mathcal C_d$ を選び，
\begin{equation}\label{ke:eq:point-anchor-target-error}
 d_\square(f_{j,t}(x),q_x)\le C\eta
\end{equation}
とできる．さらに $q_x\in\mathcal C_J$ を満たす\emph{どの}三元集合 $J$ についても，
前補題の選択に整合する $a_J(q_x)$ を取り
\begin{equation}\label{ke:eq:point-to-portal-cost}
 d_{g_t}(x,a_J(q_x))\le H
\end{equation}
とできる．
\end{lemma}

\begin{proof}
全域構成の被覆を用いて，次の三つの場合に分ける．

まず $x\in A_{t,J_0}$ かつ
$|\operatorname{supp}P^{\rm old}_{t,J_0}(x)|=3$ なら
$q_x=P^{\rm old}_{t,J_0}(x)$ とする．保持 corner 上の基底値比較により
\eqref{ke:eq:point-anchor-target-error}を得る．
$q_x$ を含む三元集合は $J_0$ に限られ，前補題で選んだ $a_{J_0}(q_x)$ は
$x$ と同じ旧 torus に属する．その径 $C/m_j\le H$ により接続できる．
この選択により，異なる endpoints が同じ anchor を持つ場合にも一つの portal を使える．

次に $x$ が pair $ab$ の保持 neck 上にあり，
$\mu(x)\in[-\delta,\delta]$ であるときは
\[
 q_x=E_{ab}(\sigma_t(\mu(x)))
\]
とする．系\ref{ke:neck-base-value-control}が target 誤差を与える．
任意の $J\supset\{a,b\}$ で
$a_J(q_x)\in\Lambda_{ab,J,t}(\mu(x))$ と選べば，
両点は同じ $P_{ab,j}$ の同一 moment 集合に属し，接続費用は $H$ 以下である．
保持 corner 上で旧 $P^{\rm old}$ の support が二元になる点もこの場合に含まれる．
実際，その全 torus は $B=m_j+O(1)$ の旧辺上にあり，
K2 の共通 plateau overlap 内の同じ neck fiber である．
Core cut により残る旧辺の interval は正確に $\mu\in[-\delta,\delta]$ である．

残る点は一つの component $a$ の core 又は変更 collar に含まれる．
$q_x=e_a$ とすれば，命題\ref{ke:global-base-parameterization}の core-image 評価が
target 誤差を与える．各 $J\ni a$ の portal を
$S_{a,j,t}\cap W_{t,J}$ に選べば，$x$ と同じ固定 component compact に属するので
接続費用は $H$ 以下である．
Anchor moment interval の外側の neck 尾部は
補題\ref{ke:component-neck-tail-depth}により同じ対数予算で扱える．

複数の場合が重なる点はどちらで扱ってもよい．この分類の役割は，
gallery の最初の square の label が元の corner の label と異なるときにも，
同じ moment 又は同じ component 内で最初の portal へ到達できることにある．
\end{proof}

\subsection{有限 gallery による全域上界}

\begin{theorem}[Actual 全曲面距離の一様上界]
\label{ke:global-ambient-upper}
任意の固定 $d\ge5$ と上の actual fibration に対し
\[
 M_d=3\binom{\min\{d,6\}}3\le60
\]
と置く．固定 stage の全条件を満たす十分小さい $t$ に対して，全 $x,y\in X_t$ で
\begin{equation}\label{ke:eq:global-ambient-upper}
 \bar d_t(x,y)
 \le d_\square(f_{j,t}(x),f_{j,t}(y))+U_{j,t},
\end{equation}
ただし
\begin{equation}\label{ke:eq:global-ambient-upper-error}
 U_{j,t}=M_dE_{j,t}
       +C(M_d+1)\eta_{j,t}
       +(M_d+1)\frac{H_{j,t}}{R_t}.
\end{equation}
特に上の許容 cut 選択と $N_j\ge m_j=6b_j$ の下で
\begin{equation}\label{ke:eq:global-upper-explicit-depth-error}
 U_{j,t}\le C_d\left\{\frac1{b_j}
       +\frac{1+\log R_t+\log(N_j+1)}{R_t}\right\}.
\end{equation}
従って既存の遅い対角選択を十分小さくすれば，$j=j(t)\to\infty$ に沿って
$U_{j(t),t}\to0$ である．同じ選択で，全 $f_t$-fibers の正規化 ambient 径も
一様に零へ収束する．
\end{theorem}

\begin{proof}
前補題の anchors $q_x,q_y$ を取る．
系\ref{topo:six-label-gallery}により，両者を結ぶ最短路を
$M\le M_d$ 本の unit-square 内線分に分ける：
\[
 z_0=q_x,\ z_1,\ldots,z_M=q_y,
 \qquad[z_{r-1},z_r]\subset Q_r\subset\mathcal C_{J_r}.
\]
その Euclidean 長さを $\ell_r$ とすれば
$\sum_r\ell_r=d_\square(q_x,q_y)$ である．
両 anchors が等しい場合は，その点を含む一つの square の定値線分を使えばよい．

第 $r$ 線分の両端に対応する portals
\[
 p_r^-=a_{J_r}(z_{r-1}),\qquad p_r^+=a_{J_r}(z_r)
\]
を取る．初めと終わりでは補題\ref{ke:point-to-corner-portal}の指定点を使う．
\eqref{ke:eq:retained-protected-estimates}から得られる actual ambient 上界と
\eqref{ke:eq:corner-portal-target-error}により
\begin{align*}
 \bar d_t(p_r^-,p_r^+)
 &\le d_\square(P^{\rm old}_{t,J_r}(p_r^-),
                         P^{\rm old}_{t,J_r}(p_r^+))+E\\
 &\le \ell_r+2C\eta+E .
\end{align*}
ここで使う経路は旧 full corner $W_{t,J_r}$ の中の経路であり，
core cut で除いた cap を通っても $X_t$ 内の正当な経路である．
保持領域 $A_{t,J_r}$ の中に経路を閉じ込める必要はない．

隣接する $J_r,J_{r+1}$ が同じなら共通 portal を使える．
異なる場合，$z_r$ の support は二元以下である．
二元の中央 interval なら両 portals は同じ $\sigma_t^{-1}(u)$ を持つので，
\eqref{ke:eq:corner-portal-junction-cost}で費用 $H$ 以下の接続を得る．
Vertex に近い端 interval と一元 support では，同じ component compact の費用を使う．
最初と最後の $x,y$ への接続も各 $H$ 以下である．
従って全ての接続を足し合わせると
\[
 \bar d_t(x,y)
 \le d_\square(q_x,q_y)+M(E+2C\eta)+(M+1)H/R_t.
\]
さらに\eqref{ke:eq:point-anchor-target-error}を両端に用いれば
\[
 d_\square(q_x,q_y)
 \le d_\square(f_{j,t}(x),f_{j,t}(y))+2C\eta .
\]
$M\le M_d$ と合わせ，定数をまとめれば
\eqref{ke:eq:global-ambient-upper}--\eqref{ke:eq:global-ambient-upper-error}を得る．
距離の上界を実際の経路長で書く場合には，各距離の infimum を任意の正の誤差で
近似する有限個の経路を選び，最後にその誤差を零へ送ればよい．

最後に全ての fixed-stage 定数を先に $C_j$ に含め，既存の parameter 閾値とともに
\[
 R_t\ge j\bigl(1+C_j+\log(N_j+1)\bigr)
\]
を課す．$m_j\to\infty$，$N_j\ge m_j$ であり，
$\log R_t/R_t\to0$ も成立するので
$E_{j(t),t},\eta_{j(t),t},H_{j(t),t}/R_t\to0$ である．
これは既存の有限個の stage 条件を強めるだけの遅い対角選択である．
全 fiber 合併径については\eqref{ke:eq:full-fiber-diameter-upper}を使う．
\end{proof}

\paragraph{平方の lift 条件との区別．}
本節で実測した局所条件は，近い旧 target 値を持つ endpoints の
actual ambient 距離の上界と，junction portals の同一 moment 接続である．
命題\ref{topo:finite-portal-criterion}の \textup{(S)} は，target 線分へ
\emph{厳密に}写る rectifiable lift を要求する．Inset と core cap の再 parameter 化の
後では，旧 corner の lift がこの条件を満たすとは限らない．
上の証明はその \textup{(S)} を引用せず，必要な endpoint 版の条件を
有限 gallery に直接適用している．\textup{(F)} に相当する actual 全 fiber 合併径は
命題\ref{ke:collapsed-tree-full-fiber-collapse}で独立に証明した．
全曲面を通る任意の経路に対する下界 \textup{(N)} は本節の結論に含めない．
